\section{Two ways prediction can miss what a decision needs}
\label{sec:mechanisms}

\subsection{An exact mean can rank actions incorrectly}
\label{sec:variance-example}

For a fixed encoder, squared prediction loss is minimized by the conditional mean
$m(c,a)=\E[Z'\mid c,a]$, where $c$ is the available history. A goal planner instead needs to compare
outcome costs. For a fixed latent goal $z_g$,
\begin{equation}
\label{eq:variance-cost}
\E\!\left[\|Z'-z_g\|^2\mid c,a\right]
=\|m(c,a)-z_g\|^2+\tr\Cov(Z'\mid c,a).
\end{equation}
The first term scores the mean; the second records the spread around it. An action-dependent spread
can change which action minimizes expected cost.

Consider the exact scalar example in Figure~\ref{fig:concepts}. The goal is zero. Action A produces
$X'=-1$ or $+1$, each with probability $1/2$; action B produces $X'=0.1$ deterministically. The two
means are $0$ and $0.1$, so their squared distances to the goal are $0$ and $0.01$: a mean-goal
planner chooses A. The correct expected costs are $1$ and $0.01$: B is better. No improved estimate
of the same mean repairs this ranking. The omitted variance is exactly the missing quantity.

\begin{figure}[t]
\centering
\paperfig{redcube_concepts_paper.pdf}{1.0}
\caption{\textbf{Correct one-step means can still give wrong decisions and wrong rollouts.}
Left: the exact mean of action A lies at the goal, although every outcome is farther from it than
action B's deterministic outcome. Right: for $X_{t+1}=X_t^2+\xi_t$, with independent fair
$\xi_t\in\{-1,+1\}$ and $X_0=0$, repeated conditional-mean prediction stays at zero while the true
two-step mean is $1$. These are analytical examples, not outputs of trained image models.}
\label{fig:concepts}
\end{figure}

\subsection{An exact one-step mean need not compose into a later mean}
\label{sec:recursive-example}

Let $X_{t+1}=X_t^2+\xi_t$, with independent fair noise $\xi_t\in\{-1,+1\}$ and $X_0=0$.
The exact conditional mean is $m(x)=x^2$, so repeated mean prediction remains zero. In the real
process, $X_1$ is $-1$ or $1$ and $X_2$ is $0$ or $2$, each fairly. Consequently,
$\E X_2=1\ne m(m(0))=0$. More generally,
\begin{equation}
\label{eq:mean-recursion}
\E[X_{t+1}]=(\E X_t)^2+\Var(X_t),
\end{equation}
while recursive mean substitution retains only the first term. This discrepancy occurs with an
exact one-step predictor and finite-support noise.

These examples concern mean-state rollout and goal evaluation. A model that predicts expected
costs or suitable value functions directly can preserve the needed expectations without generating
samples \citep{wan2019expectation,grimm2021proper}. A distributional world model offers another route:
it can retain distinct outcomes and average their costs. For the single quadratic goal cost in
Equation~\ref{eq:variance-cost}, the mean plus one scalar variance trace per action already suffices. Whether a fitted model achieves that
benefit is an empirical question.

\subsection{Non-collapse does not determine which feature is retained}
\label{sec:nuisance-example}

With finite second moments, a fixed encoder, and an unrestricted predictor, the minimum squared prediction risk is the expected
conditional latent variance:
\[
\inf_g\E\|g(C,A)-Z'\|^2=\E\tr\Cov(Z'\mid C,A).
\]
Joint representation learning changes both the target and the conditioning features in this
expression. It can make prediction easy by retaining a feature that barely changes.
Write an observation as $X_t=R(P_t,B)$,
where $P_t$ is controlled state and $B$ is an irrelevant nuisance fixed for the episode. If the
nuisance is recoverable and can be mapped to $\phi(B)\sim\Nrm(0,I)$, an encoder that returns only
$\phi(B)$ is fully non-collapsed. Copying its last latent predicts the next one with zero error.
A Gaussian marginal regularizer is also zero at the population level. Thus a global optimum can
ignore the controlled state entirely.

This is the predictable-nuisance mechanism of \citet{sobal2022slow}, specialized here to the
Gaussian-regularized objective. The construction assumes a sufficiently rich nuisance and a
realizable encoder; it is not an assertion that a finite trained network reaches that optimum.
It explains why non-collapse alone cannot choose the feature that matters. The empirical test is
whether the learned representation continues to expose cube state when nuisance and dynamics change.
Full assumptions and proofs are in Appendix~\ref{app:proofs}.
